\section{Setup and Protocol}
\label{sec:setup}

\paragraph{Data and splits.}
We use the 50,000 labeled CIFAR-10 training images (10 classes, $32\times32$
RGB). A fixed stratified 5,000-image holdout (split seed 0) is used by every
student and every teacher; the remaining 45,000 images are used for training.
Teachers and students must hold out the same images, otherwise teacher
predictions on the holdout leak into the student through the soft targets. The
final full-data models are retrained on all 50,000 images once the recipe is
fixed. No external data and no pretrained weights are used.

\paragraph{Evaluation.}
The competition metric is top-1 accuracy. The public and private leaderboards each
use 1,000 hidden images, so a public score near 90\% has a standard error of about
1.1 percentage points; we do not select models by small public differences. Unless
stated otherwise, validation numbers use EMA weights. We report raw, EMA, flip
TTA (average of the image and its horizontal mirror), and 10-view TTA. The latter,
which we call \emph{5-crop} to match the competition codebase, averages the
softmax over a center view and four $\pm2$-pixel shifts, each with and without a
horizontal flip.

\paragraph{Paired statistics.}
Model comparisons use McNemar's exact test on per-image correctness over the same
5,000 holdout images. We report the discordant counts and two-sided binomial
$p$-values. The seed-to-seed spread on this holdout is about 0.2pp on clean data
and 0.4pp on the composite proxy, so differences of that size are treated as
noise.

\paragraph{Rules.}
The official overview asks for a modified ResNet with at most 5M trainable
parameters, trained from scratch on the permitted CIFAR-10 training data, with no
pretrained weights or external datasets. The evaluation images may not be
labeled, used for training, or used for model selection. We apply the budget to
the submitted network; the teachers are auxiliary models trained from scratch on
the same 50,000 images and are not part of the submitted inference system. The
team submits a single network per submission.

\paragraph{Validation proxies.}
We build shifted validation sets from the 5,000 holdout images only. The
corruptions are Gaussian noise ($\sigma=0.05,0.10$), Gaussian blur (radius 1.0,
2.0), JPEG compression (quality 30), contrast scaling (0.7, 1.3), color
saturation (1.3), and a composite \emph{combo} (contrast 1.15, brightness 1.05,
blur 0.5, Gaussian noise $\sigma=0.02$, JPEG 75). The photographic proxies apply
per-image brightness/contrast/saturation scaling plus a small zoom and rotation:
\emph{mild} uses $\pm8\%$ brightness/contrast, $\pm12\%$ saturation, zoom 1.05 and
$\pm3^\circ$; \emph{strong} uses $\pm20\%$, $\pm30\%$, zoom 1.15 and $\pm6^\circ$.
These are label-free stress tests built from training images, not models of the
final test set.

\paragraph{Compute.}
Main runs used NYU HPC A100 nodes; ablations used A10-class cloud GPUs and a local
RTX 2070. A 4.9M student takes about 19--21s per epoch on an A10-class GPU;
three-teacher distillation takes about 59s per epoch. The longest runs (1,200
epochs on all 50k images) take several hours.
